\documentclass[10pt,a4paper]{amsart}
\usepackage{amsmath,amssymb,amsthm,graphicx}
% Retained preamble for the Stacks Project chapter snapshot in this directory.
% The chapter's own preamble.tex is a copy of the shared Stacks Project preamble; the block
% below is sliced MECHANICALLY, line for line, from it. Nothing was added, removed or
% redefined; the wrapper class, the xr-hyper/hyperref wiring and the externaldocument list
% are omitted because the assembler supplies the document class.
%
% lines taken from preamble.tex: 3, 5, 9, 11, 12, 156, 157, 158, ... (51 lines total)
%
% package coverage restored from the chapter's own preamble (2026-10-13): xy, xy-2cell

\usepackage[all]{xy}
\xyoption{2cell}
\UseAllTwocells

\newenvironment{reference}{\comment}{\endcomment}
\newenvironment{slogan}{\comment}{\endcomment}
\newenvironment{history}{\comment}{\endcomment}
\theoremstyle{plain}
\newtheorem{theorem}[subsection]{Theorem}
\newtheorem{proposition}[subsection]{Proposition}
\newtheorem{lemma}[subsection]{Lemma}
\theoremstyle{definition}
\newtheorem{definition}[subsection]{Definition}
\newtheorem{example}[subsection]{Example}
\newtheorem{exercise}[subsection]{Exercise}
\newtheorem{situation}[subsection]{Situation}
\theoremstyle{remark}
\newtheorem{remark}[subsection]{Remark}
\newtheorem{remarks}[subsection]{Remarks}
\numberwithin{equation}{subsection}
\def\lim{\mathop{\mathrm{lim}}\nolimits}
\def\colim{\mathop{\mathrm{colim}}\nolimits}
\def\Spec{\mathop{\mathrm{Spec}}}
\def\Hom{\mathop{\mathrm{Hom}}\nolimits}
\def\Ext{\mathop{\mathrm{Ext}}\nolimits}
\def\SheafHom{\mathop{\mathcal{H}\!\mathit{om}}\nolimits}
\def\SheafExt{\mathop{\mathcal{E}\!\mathit{xt}}\nolimits}
\def\Sch{\mathit{Sch}}
\def\Mor{\mathop{\mathrm{Mor}}\nolimits}
\def\Ob{\mathop{\mathrm{Ob}}\nolimits}
\def\Sh{\mathop{\mathit{Sh}}\nolimits}
\def\NL{\mathop{N\!L}\nolimits}
\def\CH{\mathop{\mathrm{CH}}\nolimits}
\def\proetale{{pro\text{-}\acute{e}tale}}
\def\etale{{\acute{e}tale}}
\def\QCoh{\mathit{QCoh}}
\def\Ker{\mathop{\mathrm{Ker}}}
\def\Im{\mathop{\mathrm{Im}}}
\def\Coker{\mathop{\mathrm{Coker}}}
\def\Coim{\mathop{\mathrm{Coim}}}
\DeclareMathSymbol{\boxtimes}{\mathbin}{AMSa}{"02}
\def\QCohstack{\mathcal{QC}\!\mathit{oh}}
\def\Cohstack{\mathcal{C}\!\mathit{oh}}
\def\Spacesstack{\mathcal{S}\!\mathit{paces}}
\def\Quotfunctor{\mathrm{Quot}}
\def\Hilbfunctor{\mathrm{Hilb}}
\def\Curvesstack{\mathcal{C}\!\mathit{urves}}
\def\Polarizedstack{\mathcal{P}\!\mathit{olarized}}
\def\Complexesstack{\mathcal{C}\!\mathit{omplexes}}
\def\Pic{\mathop{\mathrm{Pic}}\nolimits}
\def\Picardstack{\mathcal{P}\!\mathit{ic}}
\def\Picardfunctor{\mathrm{Pic}}
\def\Deformationcategory{\mathcal{D}\!\mathit{ef}}
\title{Finite type modules have a set of isomorphism classes}
\author{The Stacks Project Authors}
\date{stacks-project:examples-stacks.tex:local-snapshot-20261012; GNU FDL 1.2 or later}
\begin{document}
\maketitle
\noindent\textit{Editorial note.} Statement span is lines 449-458 of sources/stacks-examples-stacks/pinned-original-source.tex; proof span is lines 460-486, the complete author proof (25 body lines).\par\medskip
\noindent\textit{Editorial note.} Navigation stubs: 0 cited-but-not-included labels are supplied as non-load-bearing pointers with their pinned tags where the snapshot tag index has them.\par\medskip
\noindent\textit{Editorial note.} extra\_packages is empty on purpose: the retained preamble already loads xy with options and a second declaration clashes in the assembler wrapper.\par\medskip
\noindent\textit{Editorial note.} Text is retained verbatim, including the author's own typos and placeholders.\par\medskip
\noindent\textit{Editorial note.} Queue index 6 of queue\_examples-stacks.json (pb=1339, stmt 449-458, proof 460-486, halfA).\par\medskip
\noindent\textit{Editorial note.} stacks-examples-stacks wave-1 (exs1): fresh adjudication accept (2026-10-14); pool per Chief 09:4x authorization (idle-line throughput over kun's backup note, SMC precedent); examples-stacks.tex 27 labels/20 with proof judged 9 accept / 11 reject; claim\_key v2 stacks:examples-stacks.tex:<label>.\par\medskip
\noindent\textit{Allowed standard prerequisites (editorial).}\begin{itemize}
\item Explicit prerequisite boundary. The cited inputs are prerequisites; the argument in the retained span is not replaced by any cited result.
\end{itemize}\medskip
\begin{lemma}
\label{lemma-finite-type}
Let $(X, \mathcal{O}_X)$ be a ringed space.
\begin{enumerate}
\item The category of finite type $\mathcal{O}_X$-modules has a
set of isomorphism classes.
\item The category of finite type quasi-coherent
$\mathcal{O}_X$-modules has a set of isomorphism classes.
\end{enumerate}
\end{lemma}

\begin{proof}
Part (2) follows from part (1) as the category in (2) is a full subcategory
of the category in (1). Consider any open covering
$\mathcal{U} : X = \bigcup_{i \in I} U_i$. Denote $j_i : U_i \to X$
the inclusion maps. Consider any map $r : I \to \mathbf{N}$.
If $\mathcal{F}$ is an $\mathcal{O}_X$-module whose restriction to
$U_i$ is generated by at most $r(i)$ sections from $\mathcal{F}(U_i)$,
then $\mathcal{F}$ is a quotient of the sheaf
$$
\mathcal{H}_{\mathcal{U}, r} =
\bigoplus\nolimits_{i \in I} j_{i, !}\mathcal{O}_{U_i}^{\oplus r(i)}
$$
By definition, if $\mathcal{F}$ is of finite type, then there exists
some open covering with $\mathcal{U}$ whose index set is $I = X$
such that this condition is true. Hence it suffices to show that
there is a set of possible choices for $\mathcal{U}$ (obvious),
a set of possible choices for $r : I \to \mathbf{N}$ (obvious), and
a set of possible quotient modules of $\mathcal{H}_{\mathcal{U}, r}$
for each $\mathcal{U}$ and $r$. In other words, it suffices to show
that given an $\mathcal{O}_X$-module $\mathcal{H}$ there is at most
a set of isomorphism classes of quotients.
This last assertion becomes obvious
by thinking of the kernels of a quotient map
$\mathcal{H} \to \mathcal{F}$
as being parametrized by a subset of the power set of
$\prod_{U \subset X\text{ open}} \mathcal{H}(U)$.
\end{proof}

\end{document}
